\section{序論}
\label{sec:introduction}

本論文はVED$\times$IFGT枠組み\citep{VED_vol1, IFGT}の中で、概念構造と思考のダイナミクスに関するモデル内の解釈を提案する。
Master Variable Table v2.1との意図された記法の連続性は保持するが、その正確な出典・版と、型を明示した継承は未解決である。
ここで展開する枠組みをConceptual Topography（CT、概念地形）と呼ぶ。
直観的な比喩を強調するときの「概念的な地形」（conceptual terrain）も同じ対象を指す。

\subsection{概念構造の問題}
概念とは何か。
標準的な答えでは、概念を意味の領域に位置づける。
すなわち、意味を担う単位、心的表象、あるいは命題ネットワークのノードとする。
この答えは多くの目的に有用だが、\emph{動的な}問いを残す。
概念はどのように形成され、安定化し、その後の思考を制約するのか。
理解はどのように生じ、関連する定義が利用可能でも、ときに別の状況へ転移しないのはなぜか。

本論文はこれらの問いに別の方向から接近する。
概念が何を\emph{意味する}かではなく、情報場のある領域が持続し、低コストで再進入でき、将来のダイナミクスに系統的な偏りを与えるには、どのような構造条件が必要かを問う。
この捉え直しによって、問題を意味論から幾何へ移す。

\subsection{\texorpdfstring{VED$\times$IFGT}{VED x IFGT}シリーズ内の位置}
VED$\times$IFGT枠組みは二つの基礎層からなる。
Vortical Enclosure Dynamics（VED）は、差異、勾配、流れ、閉包からの構造生成を記述する\citep{VED_vol1}。
Information Field Geometry Theory（IFGT）は、$C$が1に達しない領域で準閉包として持続する、相補的な情報構造を記述する\citep{IFGT}。

本論文をこの枠組みの\emph{応用層1}に位置づけ、IFGTの場理論的語彙に認知的・社会的解釈を提案する。
新たな基礎変数は導入しない。VED$\times$IFGTの方程式は、認知現象を整理する構造クラスの候補を動機づけるものであり、確立された存在定理や分類定理ではない。

シリーズの構成は次のとおりである。
\begin{center}
\begin{tabular}{ll}
\toprule
層 & 理論 \\
\midrule
L0（生成） & VED \\
L0（構造） & IFGT \\
L1（応用） & \textbf{Conceptual Topography［本論文］} \\
L1（応用） & Consciousness Theory［repository内の関連論文］ \\
L1（応用） & Intelligence Theory［repository内の関連論文群］ \\
\bottomrule
\end{tabular}
\end{center}

\subsection{主要な主張}
本論文は局所的な定義を保持し、以下の解釈を提案する。
概念と理解の定義は、それらと併せて論じる認知的・現象的対応を確立するものではない。

\begin{enumerate}
\item \textbf{概念}は、情報ポテンシャル$\Phi(x,\tau)$の局所的安定領域であり、小さな摂動に対する安定性を伴う$\nabla\Phi(x)\approx 0$によって定義する。
\item \textbf{理解}は、概念領域を通過するコストの一時的な低下であり、正しい意味の獲得ではない。
\item \textbf{思考}は、形式的な情報流$\JI$に関連する軌道を用いて解釈する。この記述は$\nabla\Phi$による偏りを含むが、勾配だけで完全な経路を指定するものではなく、モデル内で指揮する主体を必要としない。
\item \textbf{言語トークン}は、意味内容を担わずに$\nabla\Phi$へ局所的に結合する、離散的で再アドレス可能な単位としてモデル化する。再利用による深化は仮説であり、型付けと更新に関する帰結は未解決である。
\item モデル内では、厳密な正値性$I(x,\tau)>0$を、通常の完全な概念的閉包を排するものと解釈する。さらなる分化は条件付きの提案であり、局所的な静穏と両立し、正値性だけから各点で含意されるものではない。
\item \textbf{社会的スケール}では、トークンの役割の比較が、共有地形の提案、制度をアトラクタとみる類比、集合的な維持についての局所的記述を動機づける。共通のダイナミクスや知性の一般的定義が確立されたわけではない。
\end{enumerate}

\subsection{論文の構成}
第\ref{sec:intuitive}節では概念地形の直観的な描像を示す。
第\ref{sec:formal}節ではVED$\times$IFGTの語彙から概念的現象への対応を提案する。
第\ref{sec:core}節では既存の場と結合の表式を、局所的定義および未解決の形式的範囲とともに示す。
第\ref{sec:structural}節では構造クラスの候補を整理する。
第\ref{sec:dynamics}節では沈殿、忘却、概念変化についての定性的仮説を展開する。
第\ref{sec:multiscale}節では個体から社会的地形へのスケール反転を扱う。
第\ref{sec:pathologies}節では失敗領域の提案と限界を述べる。
第\ref{sec:note}節では、成果物の帰属が未解決であることを保持しつつ、形式構造と記録されたノートの直観的語彙を結びつける。
第\ref{sec:discussion}節で含意と未解決問題を論じ、第\ref{sec:conclusion}節で結ぶ。
